%%%%%%%%%%%%%%%%%%%%%%%%%%%%%%%% 
% 该模板现在仍有缺陷：当正文内容发生跨页时不会自动封底封顶，一个可能的解决办法是写完报告后手动加封底线、封顶线。
% 你也可以自行改进该模板，增加自动封顶封底功能。
%%%%%%%%%%%%%%%%%%%%%%%%%%%%%%%%

\documentclass[a4paper, zihao=5, UTF8, AutoFakeBold]{ctexart}

%% ===== 参考文献 =====
\usepackage[backend=biber, style=gb7714-2015]{biblatex}
\addbibresource{references.bib}

%% ===== 排版宏包 =====
\usepackage[most]{tcolorbox}
\usepackage{tabularray}
\usepackage{geometry}
\usepackage{amsmath}
\usepackage{amssymb}
\usepackage{graphicx}
\usepackage{float}
\usepackage{xcolor}
\usepackage{hyperref}

%% ===== 页面设置 =====
\geometry{top=2.5cm, bottom=2cm, left=2.5cm, right=2.5cm}

%% ===== 字体命令 =====
\newcommand{\bfsong}{\songti\zihao{-4}\bfseries}

%% ===== 内容区左右内边距 =====
\newlength{\cpad}
\setlength{\cpad}{5pt}

%% ===== 大框样式 =====
\tcbset
{
    reportbox/.style=
    {
        breakable, enhanced jigsaw,
        colback  = white,
        colframe = black,
        boxrule  = 0.6pt,
        arc      = 0pt,
        left = 0pt, right = 0pt,
        top  = 0pt, bottom = 0pt,
        boxsep   = 0pt,
    },
    segmentation style=
    {
        solid,
        draw=black,
        line width=0.6pt
    }
}

%% ===== 节标题命令 =====
\newcommand{\secheader}[1]
{%

    \noindent\hspace{\cpad}{\songti\zihao{-4}\bfseries #1}%

}

%% ===== 正文内容区环境（支持跨页）=====
\newenvironment{contentarea}
{%
    \par\vspace{4pt}%
    \leftskip=\cpad\rightskip=\cpad%
    \setlength{\parindent}{2em}%
    \par\vspace{4pt}%
}

% ============================================================

\begin{document}
\pagestyle{empty}

%% ===== 标题 =====
\begin{center}
    {\songti\zihao{-2}\bfseries 大学物理实验三设计性实验调研报告}
\end{center}

%% ===== 大框开始 =====
\begin{tcolorbox}[reportbox]


%% ---- 信息填写区（tblr 无外框，tcolorbox 提供外边框）----
\begin{tblr}
    {
        width    = \linewidth,
        colspec  = { Q[c,m] X[c,m] Q[c,m] X[c,m] Q[c,m] X[c,m] },
        cell{2}{2} = {c=5}{c,m},
        hline{2,3} = {0.6pt},
        vline{2,3,4,5,6} = {0.6pt},
        row{1,2} = {ht=0.75cm},
    }
{\bfsong 课程编号} & 1800460001 & {\bfsong 组号} & 3 & {\bfsong 姓名} & 请在这里填写姓名 \\
{\bfsong 实验名称} & 请在此处填写实验名称 &&&& \\
\end{tblr}



%% ---- 一、选题意义 ----

\par\vspace{4pt}

\secheader{一、选题意义}

\tcbline

\begin{contentarea}
% 在此处填写选题意义
% 要阐明为什么要研究这个题目、研究它有什么价值、能解决什么问题

\noindent 示例内容\\示例内容\\示例内容\\示例内容\\示例内容\\示例内容\\示例内容\\示例内容\\示例内容\\示例内容\\
示例内容\\示例内容\\示例内容\\示例内容\\示例内容\\示例内容\\示例内容\\示例内容\\示例内容\\示例内容\\

\end{contentarea}

\tcbline

%% ---- 二、文献归纳 ----
\secheader{二、文献归纳}

\tcbline

\begin{contentarea}
% 在此处填写文献归纳，用 \cite{key} 在正文中引用
% 对他人的实验思路、方法、成果进行归纳整理

\noindent 示例内容\\示例内容\\示例内容\\示例内容\\示例内容\\示例内容\\示例内容\\示例内容\\示例内容\\示例内容\\
示例内容\\示例内容\\示例内容\\示例内容\\示例内容\\示例内容\\示例内容\\示例内容\\示例内容\\示例内容\\

% 参考文献列表（自动生成，勿手动修改）
\printbibliography[heading=none]
\end{contentarea}

\tcbline

%% ---- 三、预期目标 ----
\secheader{三、预期目标}

\tcbline

\begin{contentarea}
% 在此处填写预期目标
% 简述拟采用何种方法，通过实验解决什么问题，达到什么目标

\noindent 示例内容\\示例内容\\示例内容\\示例内容\\示例内容\\示例内容\\示例内容\\示例内容\\示例内容\\示例内容\\
示例内容\\示例内容\\示例内容\\示例内容\\示例内容\\示例内容\\示例内容\\示例内容\\示例内容\\示例内容\\

\end{contentarea}

\tcbline

%% ---- 四、指导教师意见 ----
\secheader{四、指导教师意见}

\tcbline

\vspace{3cm} % 教师意见留白区，如需调整高度请修改此数值

%% ---- 成绩/签章行 ----
\begin{tblr}
    {
        width    = \linewidth,
        colspec  = { Q[c,m] X[c,m] Q[c,m] X[c,m] X[c,m] X[c,m] },
        hline{1} = {0.6pt},
        row{1} = {ht=0.75cm},
    }
{成绩} &  & {签章:} &  & {年\qquad 月\qquad 日}\\
\end{tblr}

\end{tcolorbox}

% \vspace{6pt}
% {\zihao{-5}\kaishu 附注：本调研报告成绩占期末考试成绩的30\%}

\end{document}